\documentclass[11pt]{article}

\usepackage[a4paper,margin=1in]{geometry}
\usepackage{amsmath,amssymb,amsthm,mathtools,bm}
\usepackage{enumitem}
\usepackage{booktabs,tabularx,array}
\usepackage{microtype}
\usepackage[hidelinks]{hyperref}
\usepackage{graphicx}
\usepackage{tikz}
\usetikzlibrary{arrows.meta,positioning,shapes.geometric}

\newtheorem{assumption}{Assumption}[section]
\newtheorem{definition}[assumption]{Definition}
\newtheorem{proposition}[assumption]{Proposition}
\newtheorem{theorem}[assumption]{Theorem}
\newtheorem{lemma}[assumption]{Lemma}
\newtheorem{corollary}[assumption]{Corollary}
\theoremstyle{remark}
\newtheorem{remark}[assumption]{Remark}
\newtheorem{example}[assumption]{Example}

\newcommand{\E}{\mathbb E}
\newcommand{\Pp}{\mathbb P}
\newcommand{\R}{\mathbb R}
\newcommand{\Var}{\operatorname{Var}}
\newcommand{\Cov}{\operatorname{Cov}}
\newcommand{\diag}{\operatorname{diag}}
\newcommand{\op}{\mathrm{o}_{\Pp}}
\newcommand{\Op}{\mathrm{O}_{\Pp}}
\newcommand{\cf}{\mathrm{cf}}
\newcommand{\fact}{\mathrm{fact}}
\newcommand{\cW}{\mathcal W}
\newcommand{\cT}{\mathcal T}
\newcommand{\cC}{\mathcal C}
\newcommand{\cG}{\mathcal G}
\newcommand{\cU}{\mathcal U}
\newcommand{\norm}[1]{\left\lVert #1\right\rVert}
\newcommand{\abs}[1]{\left\lvert #1\right\rvert}
\newcommand{\an}{a_{N,n}}
\newcommand{\bn}{b_{N,n}}
\newcommand{\LamN}{\Lambda_N}
\newcommand{\thetahat}{\widehat\theta_{N,n}}
\newcommand{\sigmahat}{\widehat\sigma_{N,n}}
\newcommand{\vhat}{\widehat\vartheta_{N,n}}
\newcommand{\vzero}{\vartheta_0}
\newcommand{\Jhat}{\widehat J_{B,N,n}}
\newcommand{\Dmat}{D_{N,n}}
\newcommand{\Amat}{A_{N,n}}
\newcommand{\BL}{\mathrm{BL}_1}
\newcommand{\Ph}{\Phi}
\newcommand{\ph}{\phi}

\title{Work Package 5: Feasible Multiplier Bootstrap and Simultaneous Inference for Linear SPDE Counterfactuals\\
\large Mode-block scores, mixed-rate confidence bands, contact sets, and numerical-bias control}
\author{}
\date{7 July 2026}

\begin{document}
\maketitle

\begin{abstract}
This note closes the uniform-inference stage of the operator-intervention SPDE programme. Work Package 3 established a mixed-rate locally asymptotically normal experiment for the exact spectral likelihood, and Work Package 4 propagated that experiment to compact surfaces of counterfactual CDFs, factual laws, distributional effects, quantiles, and tails. The present package constructs feasible conditional critical values. The independent resampling unit is a complete modal score block, not an individual time increment. A Gaussian mode-block multiplier yields, conditionally on the data, a two-dimensional normal perturbation with estimated inverse-information covariance. Because every Work Package 4 limit is finite-rank, the score multiplier is conditionally equivalent to a direct Gaussian draw from that covariance; no high-dimensional empirical-process approximation is required. We prove conditional weak convergence jointly for the fast fixed-intervention and slow factual/effect processes, construct raw and floor-studentised simultaneous bands, and handle derivative zeros without dividing by vanishing standard errors. Equality, integrated, stochastic-dominance, tail-effect, and quantile-effect tests are developed. One-sided composite nulls use an estimated contact set with a generalized-moment-selection threshold. We also state grid and numerical-bias conditions, including bias-enlarged bands that do not require undersmoothing. An exact-OU Monte Carlo experiment across three $(N,n,K)$ designs produces simultaneous coverage close to the nominal level. The programme audit places the project at the end of Stage C: Work Packages 1--5 are complete, while operator sensitivity, total approximation error, non-identification, full validation, and manuscript reconstruction remain.
\end{abstract}

\section{Location in the research programme}

The intended paper establishes the chain
\[
\text{intervention semantics}
\Longrightarrow
\text{identification}
\Longrightarrow
\text{efficient estimation}
\Longrightarrow
\text{uniform target processes}
\Longrightarrow
\text{feasible simultaneous inference}
\Longrightarrow
\text{total numerical--statistical error}
\Longrightarrow
\text{sharp non-identification}.
\]
The current status is recorded in Table~\ref{tab:programme}.

\begin{table}[ht]
\centering
\small
\begin{tabularx}{\textwidth}{@{}lXl@{}}
\toprule
Package & Scientific role & Status\\
\midrule
WP1 & Galerkin structural intervention semantics & Complete\\
WP2 & Target-specific identification & Complete\\
WP3 & Exact likelihood, mixed-rate LAN, efficient MLE & Theorem package complete; final appendix expansion remains\\
WP4 & Uniform counterfactual plug-in processes & Complete\\
WP5 & Feasible multiplier bootstrap and simultaneous bands & Completed here\\
WP6 & Operator sensitivity and causal response geometry & Pending\\
WP7 & Galerkin, temporal, and total inferential error & Pending\\
WP8 & Sharp non-identification boundary & Pending\\
WP9 & Full Monte Carlo and reproducibility suite & Pending; partial diagnostics exist\\
WP10 & Manuscript reconstruction & Pending\\
\bottomrule
\end{tabularx}
\caption{Position of WP5 in the full research programme.}
\label{tab:programme}
\end{table}

Thus five of ten packages are complete after this note. Decision Gates 1--5 are passed at theorem-package level. Gate 6, authorising full-field inferential claims, remains pending until WP7. The completed dependency path is
\[
 \mathrm{WP1}\to\mathrm{WP2}\to\mathrm{WP3}\to\mathrm{WP4}\to\mathrm{WP5}.
\]
WP6 was originally scheduled in the foundation stage but is not a dependency of WP5. It must now precede WP7.

\section{Mixed-rate statistical input}

Let $\vartheta=(\theta,\sigma)$ and let $S_{k,N,n}(\vartheta)\in\R^2$ be the exact conditional score contributed by the complete observed path of mode $k$:
\begin{equation}\label{eq:mode-score}
 S_{k,N,n}(\vartheta)
 :=
 \sum_{j=0}^{n-1}s_{k,j,n}(\vartheta).
\end{equation}
The modes are independent, while temporal dependence remains inside each block. Define
\begin{equation}\label{eq:scales}
 \an:=\sqrt{\LamN},
 \qquad
 \LamN:=\sum_{k=1}^N\lambda_k,
 \qquad
 \bn:=\sqrt{Nn},
 \qquad
 \Amat:=\diag(\an,\bn),
 \qquad
 \Dmat:=\Amat^{-1}.
\end{equation}

Work Package 3 proves, under its exact-likelihood regime,
\begin{equation}\label{eq:MLE-linear}
 \Amat(\vhat-\vzero)
 =
 J(\vzero)^{-1}
 \Dmat\sum_{k=1}^N S_{k,N,n}(\vzero)
 +\op(1),
\end{equation}
where
\begin{equation}\label{eq:J}
 J(\vzero)
 =
 \begin{pmatrix}
 T/(2\theta_0)&0\\
 0&2/\sigma_0^2
 \end{pmatrix}.
\end{equation}
In particular,
\begin{equation}\label{eq:mixed-limit}
 \Amat(\vhat-\vzero)
 \Rightarrow
 H\sim\mathcal N_2(0,J(\vzero)^{-1}).
\end{equation}

\begin{assumption}[WP3 regime]\label{ass:wp3}
The exact spectral experiment satisfies
\[
 N\to\infty,
 \qquad
 \lambda_N\Delta_n\to0,
 \qquad
 \frac{\lambda_N}{\LamN}\to0,
 \qquad
 \frac{\an}{\bn}\to0,
\]
together with the compact-parameter, known-spectrum, and negligible-initial-energy assumptions of Work Package 3. The MLE is interior with probability tending to one and satisfies \eqref{eq:MLE-linear}.
\end{assumption}

\begin{lemma}[Mode-block no-dominance]\label{lem:block}
Under Assumption~\ref{ass:wp3}, for some constant $C<\infty$,
\begin{equation}\label{eq:block-fourth}
 \sum_{k=1}^N
 \E_{\vzero}
 \norm{\Dmat S_{k,N,n}(\vzero)}^4
 \le
 C\left(
 \frac{\lambda_N}{\LamN}
 +\frac1N
 +(\lambda_N\Delta_n)^2
 \right)
 +o(1).
\end{equation}
Consequently,
\begin{align}
 \max_{k\le N}
 \norm{\Dmat S_{k,N,n}(\vzero)}
 &\overset{\Pp}{\longrightarrow}0,\label{eq:max-block}\\
 \sum_{k=1}^N
 \Dmat S_{k,N,n}(\vzero)
 S_{k,N,n}(\vzero)^\top\Dmat
 &\overset{\Pp}{\longrightarrow}J(\vzero).\label{eq:opg-true}
\end{align}
\end{lemma}

\begin{proof}[Proof structure]
Each modal score is a finite sum of Gaussian martingale chaoses of orders one and two. The drift-coordinate variance is asymptotic to $T\lambda_k/(2\theta_0)$, and its fourth moment is bounded by a constant multiple of $\lambda_k^2$. After scaling by $\LamN^{-2}$ and summing,
\[
 \frac{\sum_{k\le N}\lambda_k^2}{\LamN^2}
 \le
 \frac{\lambda_N}{\LamN}\to0.
\]
The scale-coordinate variance is $2n/\sigma_0^2$ and its fourth moment is $O(n^2)$, giving $O(Nn^2/(Nn)^2)=O(N^{-1})$. Mixed terms are controlled by Cauchy--Schwarz and the asymptotic score orthogonality supplied by $\lambda_N\Delta_n\to0$. This yields \eqref{eq:block-fourth}. Markov's inequality gives \eqref{eq:max-block}. Independence across modes, the fourth-moment bound, and the information calculation from Work Package 3 give \eqref{eq:opg-true}.
\end{proof}

\section{Feasible mode-block multiplier}

Define the estimated scaled modal score blocks
\begin{equation}\label{eq:Bhat}
 \widehat B_{k,N,n}
 :=
 \Dmat S_{k,N,n}(\vhat),
\end{equation}
and the block outer-product information estimator
\begin{equation}\label{eq:Jhat}
 \Jhat
 :=
 \sum_{k=1}^N
 \widehat B_{k,N,n}\widehat B_{k,N,n}^\top.
\end{equation}

\begin{lemma}[Feasible information consistency]\label{lem:Jhat}
Under Assumption~\ref{ass:wp3},
\begin{equation}\label{eq:Jhat-cons}
 \Jhat\overset{\Pp}{\longrightarrow}J(\vzero).
\end{equation}
Moreover,
\[
 \max_{k\le N}\norm{\widehat B_{k,N,n}}
 \overset{\Pp}{\longrightarrow}0.
\]
\end{lemma}

\begin{proof}
Taylor-expand each modal score between $\vzero$ and $\vhat$. Consistency and the mixed local rate imply that the sum of the scaled plug-in differences is $o_{\Pp}(1)$ in the outer-product norm. Uniform Gaussian moment bounds control the quadratic remainder. Combine this with Lemma~\ref{lem:block}.
\end{proof}

Let $\xi_1,\ldots,\xi_N$ be iid mean-zero, variance-one multipliers independent of the data. Define
\begin{equation}\label{eq:bootstrap-score}
 \Delta_{N,n}^*
 :=
 \sum_{k=1}^N\xi_k\widehat B_{k,N,n},
 \qquad
 H_{N,n}^*
 :=
 \Jhat^{-1}\Delta_{N,n}^*.
\end{equation}
If necessary, replace $\Jhat^{-1}$ by an eigenvalue-floored inverse with a floor converging to zero.

\begin{proposition}[Gaussian multiplier shortcut]\label{prop:shortcut}
If $\xi_k\sim\mathcal N(0,1)$, then conditionally on the data,
\begin{equation}\label{eq:conditional-normal}
 \Delta_{N,n}^*\sim\mathcal N_2(0,\Jhat),
 \qquad
 H_{N,n}^*\sim\mathcal N_2(0,\Jhat^{-1}).
\end{equation}
Therefore the mode-block score multiplier is conditionally identical in law to a direct draw
\[
 H_{N,n}^*=\Jhat^{-1/2}Z^*,
 \qquad
 Z^*\sim\mathcal N_2(0,I_2).
\]
\end{proposition}

\begin{proof}
A linear combination of independent Gaussian multipliers is Gaussian, with conditional covariance
\[
 \sum_{k=1}^N\widehat B_{k,N,n}\widehat B_{k,N,n}^\top=\Jhat.
\]
Premultiplication by $\Jhat^{-1}$ gives \eqref{eq:conditional-normal}.
\end{proof}

\begin{remark}[Plug-in LAN covariance]
In the positive core, one may replace $\Jhat$ by
\[
 J(\vhat)
 =
 \diag\left(\frac{T}{2\thetahat},\frac{2}{\sigmahat^2}\right).
\]
The resulting direct Gaussian multiplier is first-order equivalent and is more stable when $N$ is small. The block outer-product version remains the structural bootstrap and is the preferred template once additional nuisance coordinates are introduced.
\end{remark}

\begin{theorem}[Conditional estimator bootstrap]\label{thm:parameter-bootstrap}
Under Assumption~\ref{ass:wp3},
\begin{equation}\label{eq:H-bootstrap}
 H_{N,n}^*
 \Rightarrow_{\Pp}
 H\sim\mathcal N_2(0,J(\vzero)^{-1}).
\end{equation}
Equivalently,
\begin{equation}\label{eq:BL}
 \sup_{f\in\BL}
 \left|
 \E^*f(H_{N,n}^*)-\E f(H)
 \right|
 \overset{\Pp}{\longrightarrow}0.
\end{equation}
The result holds for Gaussian multipliers. It also holds for Rademacher multipliers under the conditional Lindeberg condition supplied by Lemmas~\ref{lem:block}--\ref{lem:Jhat}.
\end{theorem}

\begin{proof}
For Gaussian multipliers, Proposition~\ref{prop:shortcut} and Lemma~\ref{lem:Jhat} reduce the claim to continuity of the Gaussian law in its positive-definite covariance matrix. For Rademacher multipliers, condition on the data and apply the triangular-array Lindeberg--Feller theorem to $\sum_k\xi_k\widehat B_{k,N,n}$; the conditional covariance converges to $J(\vzero)$ and the largest summand vanishes.
\end{proof}

\section{Bootstrap target processes}

Use the notation of Work Package 4. A fixed absolute intervention $do(\theta=a)$ gives
\[
 F_{a,w,z}^{\cf}(\sigma)
 =
 \Ph\left(
 \frac{z-m_a(w)}
 {\sigma\sqrt{V_a(w)}}
 \right).
\]
Let
\begin{equation}\label{eq:dsigma}
 d_\sigma^{\cf}(a,w,z)
 :=
 -\frac{c_{a,w,z}(\sigma_0)\ph(c_{a,w,z}(\sigma_0))}
 {\sigma_0}.
\end{equation}
Let $d_\theta^{\fact}(w,z)$ be the factual CDF derivative defined in Work Package 4. Write $H_{N,n,\theta}^*$ and $H_{N,n,\sigma}^*$ for the coordinates of $H_{N,n}^*$.

Define
\begin{align}
 \mathbb Z_{\cf,N,n}^*(a,w,z)
 &:=
 \widehat d_\sigma^{\cf}(a,w,z)
 H_{N,n,\sigma}^*,\label{eq:Zcf}\\
 \mathbb Z_{\fact,N,n}^*(w,z)
 &:=
 \widehat d_\theta^{\fact}(w,z)
 H_{N,n,\theta}^*,\label{eq:Zfact}\\
 \mathbb Z_{\Delta,N,n}^*(a,w,z)
 &:=
 -\widehat d_\theta^{\fact}(w,z)
 H_{N,n,\theta}^*.\label{eq:ZDelta}
\end{align}

For $z_p=\Ph^{-1}(p)$, define
\begin{align}
 \mathbb Z_{Q,\cf,N,n}^*(a,w,p)
 &:=
 \sqrt{\widehat V_a(w)}\,z_p H_{N,n,\sigma}^*,\label{eq:ZQcf}\\
 \mathbb Z_{Q,\fact,N,n}^*(w,p)
 &:=
 \left[
 \widehat{\dot m}_{\thetahat}(w)
 +\sigmahat
 \frac{\widehat{\dot V}_{\thetahat}(w)}
 {2\sqrt{\widehat V_{\thetahat}(w)}}z_p
 \right]H_{N,n,\theta}^*.\label{eq:ZQfact}
\end{align}
Tail bootstrap processes are the negatives of their CDF counterparts.

\begin{assumption}[Uniform derivative consistency]\label{ass:derivative}
Every derivative surface appearing in \eqref{eq:Zcf}--\eqref{eq:ZQfact} is continuous on its compact index set and its plug-in estimator converges uniformly in probability to the true derivative surface.
\end{assumption}

\begin{theorem}[Joint conditional target-process convergence]\label{thm:process-bootstrap}
Under Assumptions~\ref{ass:wp3} and \ref{ass:derivative},
\[
 \left(
 \mathbb Z_{\cf,N,n}^*,
 \mathbb Z_{\fact,N,n}^*,
 \mathbb Z_{\Delta,N,n}^*,
 \mathbb Z_{Q,\cf,N,n}^*,
 \mathbb Z_{Q,\fact,N,n}^*
 \right)
 \Rightarrow_{\Pp}
 \left(
 \mathbb G_{\cf},
 \mathbb G_{\fact},
 \mathbb G_\Delta,
 \mathbb Q_{\cf},
 \mathbb Q_{\fact}
 \right)
\]
in the product of the continuous-function spaces used in Work Package 4.
\end{theorem}

\begin{proof}
The map sending $(h_\theta,h_\sigma)$ and the derivative surfaces to the five process coordinates is continuous from a finite-dimensional Euclidean space times a product of uniform function spaces into the target product space. Apply Theorem~\ref{thm:parameter-bootstrap}, Assumption~\ref{ass:derivative}, and the conditional continuous mapping theorem.
\end{proof}

\begin{remark}[Finite-rank structure]
No entropy integral or high-dimensional Gaussian comparison is needed. The stochastic dimension is two. The index sets affect only deterministic continuous loading surfaces multiplying the two Gaussian coordinates.
\end{remark}

\section{Simultaneous confidence bands}

Let $T_0\in C(\cT)$ be one target surface, let $\widehat T$ be its estimator, let $r_{N,n}$ be its WP4 rate, and let $\mathbb Z_{N,n}^*$ be its bootstrap process. Assume
\begin{equation}\label{eq:target-clt}
 r_{N,n}(\widehat T-T_0)\Rightarrow\mathbb G
 \quad\text{in }C(\cT),
 \qquad
 \mathbb Z_{N,n}^*\Rightarrow_{\Pp}\mathbb G.
\end{equation}

\subsection{Raw sup-norm bands}

Let
\begin{equation}\label{eq:craw}
 c_{1-\alpha}^{*,\mathrm{raw}}
 :=
 q_{1-\alpha}^*
 \left(\norm{\mathbb Z_{N,n}^*}_\infty\right).
\end{equation}
Define
\begin{equation}\label{eq:raw-band}
 \cC_{N,n}^{\mathrm{raw}}(\tau)
 :=
 \left[
 \widehat T(\tau)-\frac{c_{1-\alpha}^{*,\mathrm{raw}}}{r_{N,n}},
 \widehat T(\tau)+\frac{c_{1-\alpha}^{*,\mathrm{raw}}}{r_{N,n}}
 \right].
\end{equation}

\begin{theorem}[Raw simultaneous coverage]\label{thm:raw-band}
Suppose the distribution function of $\norm{\mathbb G}_\infty$ is continuous at its $(1-\alpha)$ quantile. Then
\begin{equation}\label{eq:raw-coverage}
 \Pp\left(
 T_0(\tau)\in\cC_{N,n}^{\mathrm{raw}}(\tau)
 \text{ for every }\tau\in\cT
 \right)
 \longrightarrow1-\alpha.
\end{equation}
For CDF and tail targets, intersecting \eqref{eq:raw-band} with $[0,1]$ preserves coverage.
\end{theorem}

\begin{proof}
Conditional weak convergence and quantile consistency imply
$c_{1-\alpha}^{*,\mathrm{raw}}\to_{\Pp}c_{1-\alpha}$, the corresponding quantile of $\norm{\mathbb G}_\infty$. Equation \eqref{eq:raw-coverage} is then the continuous-mapping consequence of \eqref{eq:target-clt}. Clipping cannot remove a true probability from an interval that already contains it.
\end{proof}

\subsection{Floor-studentised bands}

Let
\begin{equation}\label{eq:s}
 s(\tau):=\sqrt{\Var(\mathbb G(\tau))},
 \qquad
 \widehat s(\tau)\overset{\Pp}{\longrightarrow}s(\tau)
 \quad\text{uniformly}.
\end{equation}
Ordinary studentisation fails at derivative zeros. Define
\begin{equation}\label{eq:floor}
 \widehat s_\kappa(\tau)
 :=
 \max\left\{
 \widehat s(\tau),
 \kappa_{N,n}\norm{\widehat s}_\infty
 \right\}.
\end{equation}

\begin{assumption}[Floor rate]\label{ass:floor}
For a fast family with rate $\bn$,
\[
 \kappa_{N,n}\downarrow0,
 \qquad
 \bn\kappa_{N,n}\to\infty.
\]
For a slow family with rate $\an$,
\[
 \kappa_{N,n}\downarrow0,
 \qquad
 \an\kappa_{N,n}\to\infty,
 \qquad
 \frac{\an}{\bn\kappa_{N,n}}\to0.
\]
The target map has a uniform quadratic Taylor remainder.
\end{assumption}

Let
\begin{equation}\label{eq:cst}
 c_{1-\alpha}^{*,\mathrm{st}}
 :=
 q_{1-\alpha}^*
 \left(
 \sup_{\tau\in\cT}
 \frac{\abs{\mathbb Z_{N,n}^*(\tau)}}
 {\widehat s_\kappa(\tau)}
 \right),
\end{equation}
and define
\begin{equation}\label{eq:st-band}
 \cC_{N,n}^{\mathrm{st}}(\tau)
 :=
 \left[
 \widehat T(\tau)
 -
 \frac{c_{1-\alpha}^{*,\mathrm{st}}\widehat s_\kappa(\tau)}
 {r_{N,n}},
 \widehat T(\tau)
 +
 \frac{c_{1-\alpha}^{*,\mathrm{st}}\widehat s_\kappa(\tau)}
 {r_{N,n}}
 \right].
\end{equation}

\begin{theorem}[Derivative-zero robust simultaneous coverage]\label{thm:st-band}
Suppose $\norm{s}_\infty>0$, Assumption~\ref{ass:floor} holds, and the distribution of the limiting floor-studentised supremum is continuous at its $(1-\alpha)$ quantile. Then
\[
 \Pp\left(
 T_0(\tau)\in\cC_{N,n}^{\mathrm{st}}(\tau)
 \text{ for every }\tau\in\cT
 \right)
 \longrightarrow1-\alpha.
\]
At points where $s(\tau)>0$, the floor is asymptotically inactive. At points where $s(\tau)=0$, the band half-width dominates the quadratic Taylor remainder and, for a slow family, the omitted fast-coordinate contribution.
\end{theorem}

\begin{proof}[Proof structure]
On $\{s\ge\varepsilon\}$, uniform consistency makes $\widehat s_\kappa/\widehat s\to1$. On $\{s<\varepsilon\}$, first-order stochastic variation is uniformly small. The fast condition $\bn\kappa_{N,n}\to\infty$ makes a $O_{\Pp}(\bn^{-2})$ quadratic remainder negligible relative to $\kappa_{N,n}/\bn$. For a slow family, $\an\kappa_{N,n}\to\infty$ handles the $O_{\Pp}(\an^{-2})$ slow quadratic term and $\an/(\bn\kappa_{N,n})\to0$ handles the unnormalised $O_{\Pp}(\bn^{-1})$ fast nuisance term. Letting $\varepsilon\downarrow0$ and applying conditional quantile consistency yields the result.
\end{proof}

\begin{remark}[Gaussian median]
For a zero-mean fixed-intervention Gaussian law, the median is $m_a(w)$ and is exactly independent of $\sigma$. Its fast derivative vanishes identically. A separate inferential band is unnecessary at that index; the floor construction is retained so that an automated surface-level procedure remains well-defined.
\end{remark}

\subsection{Joint mixed-rate bands}

Let $\mathfrak F$ be a finite family of target surfaces, with family-specific rates $r_f$ and floor standard errors $\widehat s_{f,\kappa}$. Define
\begin{equation}\label{eq:joint-stat}
 M_{N,n}^*
 :=
 \max_{f\in\mathfrak F}
 \sup_{\tau\in\cT_f}
 \frac{\abs{\mathbb Z_{f,N,n}^*(\tau)}}
 {\widehat s_{f,\kappa}(\tau)}.
\end{equation}

\begin{corollary}[Joint fast--slow coverage]\label{cor:joint}
One conditional quantile of \eqref{eq:joint-stat}, applied with the family-specific widths
\[
 \frac{c_{1-\alpha}^{*,\mathrm{joint}}
 \widehat s_{f,\kappa}(\tau)}{r_f},
\]
gives asymptotic simultaneous coverage $1-\alpha$ across all target families in $\mathfrak F$. The common draw $H_{N,n}^*$ preserves finite-sample correlation between the estimated drift and scale coordinates.
\end{corollary}

\section{Equality and integrated tests}

Let $G_0\in C(\cT_0)$ be an effect surface estimated by $\widehat G$ at rate $r_{N,n}$, with bootstrap process $\mathbb Z_G^*$.

For
\[
 H_0^=:\quad G_0(\tau)=0
 \quad\forall\tau\in\cT_0,
\]
define
\begin{align}
 S_{N,n}^{\mathrm{KS}}
 &:=
 r_{N,n}\norm{\widehat G}_\infty,\label{eq:KS}\\
 S_{N,n}^{\mathrm{CvM}}
 &:=
 r_{N,n}^2
 \int_{\cT_0}\widehat G(\tau)^2\,d\mu(\tau),\label{eq:CvM}
\end{align}
where $\mu$ is a finite measure. Their bootstrap analogues are
\begin{align}
 S_{N,n}^{*,\mathrm{KS}}
 &:=
 \norm{\mathbb Z_G^*}_\infty,\\
 S_{N,n}^{*,\mathrm{CvM}}
 &:=
 \int_{\cT_0}\mathbb Z_G^*(\tau)^2\,d\mu(\tau).
\end{align}

\begin{theorem}[Equality-test validity]\label{thm:equality}
Under $H_0^=$ and the process-bootstrap assumptions, bootstrap critical values from the preceding statistics yield asymptotic level $\alpha$ whenever the corresponding limit distribution is continuous at the critical quantile. Under every fixed alternative with $G_0\not\equiv0$, both tests are consistent.
\end{theorem}

\begin{proof}
Under the null, apply the continuous mapping theorem to the sup norm and the continuous quadratic functional. Under a fixed alternative, $r_{N,n}\norm{G_0}_\infty\to\infty$, and the integrated statistic diverges whenever $G_0$ is non-zero on a set of positive $\mu$ measure.
\end{proof}

These tests apply to policy-to-policy CDF contrasts, factual-versus-counterfactual distributional effects, tail effects, and quantile effects.

\section{One-sided hypotheses and estimated contact sets}

Consider
\begin{equation}\label{eq:null-one}
 H_0^\le:\quad G_0(\tau)\le0
 \quad\forall\tau\in\cT_0.
\end{equation}
Define the studentised statistic
\begin{equation}\label{eq:one-stat}
 S_{N,n}^\le
 :=
 \sup_{\tau\in\cT_0}
 \left[
 \frac{r_{N,n}\widehat G(\tau)}
 {\widehat s_\kappa(\tau)}
 \right]_+.
\end{equation}
The population contact set is
\[
 \cC_0:=\{\tau\in\cT_0:G_0(\tau)=0\}.
\]

Let $\eta_{N,n}\to\infty$ and $\eta_{N,n}/r_{N,n}\to0$. Define
\begin{equation}\label{eq:Chat}
 \widehat\cC_{N,n}
 :=
 \left\{
 \tau\in\cT_0:
 \frac{r_{N,n}\widehat G(\tau)}
 {\widehat s_\kappa(\tau)}
 \ge-\eta_{N,n}
 \right\}.
\end{equation}
If $\widehat\cC_{N,n}=\varnothing$, the procedure does not reject. Otherwise define
\begin{equation}\label{eq:one-bootstrap}
 S_{N,n}^{*,\le}
 :=
 \sup_{\tau\in\widehat\cC_{N,n}}
 \left[
 \frac{\mathbb Z_G^*(\tau)}
 {\widehat s_\kappa(\tau)}
 \right]_+.
\end{equation}

\begin{assumption}[Contact-set separation]\label{ass:contact}
For every open neighbourhood $\cU$ of $\cC_0$, there exists $c(\cU)>0$ such that
\[
 \sup_{\tau\in\cT_0\setminus\cU}G_0(\tau)
 \le-c(\cU).
\]
At least one regular contact point has positive first-order standard deviation.
\end{assumption}

\begin{theorem}[Feasible one-sided inference]\label{thm:contact}
Under \eqref{eq:null-one}, Assumption~\ref{ass:contact}, and the preceding bootstrap conditions:
\begin{enumerate}[label=(\roman*)]
 \item $\cC_0\subseteq\widehat\cC_{N,n}$ with probability tending to one;
 \item for every fixed neighbourhood $\cU$ of $\cC_0$, $\widehat\cC_{N,n}\subseteq\cU$ with probability tending to one;
 \item the test rejecting when $S_{N,n}^\le$ exceeds the conditional $(1-\alpha)$ quantile of \eqref{eq:one-bootstrap} has asymptotic size at most $\alpha$;
 \item at regular boundary nulls, the asymptotic size is $\alpha$;
 \item if $\sup_{\tau\in\cT_0}G_0(\tau)>0$, the test is consistent.
\end{enumerate}
\end{theorem}

\begin{proof}[Proof structure]
At a contact point, the normalised estimator is $O_{\Pp}(1)$, while $\eta_{N,n}\to\infty$, proving inclusion. Away from contact, the deterministic term is at most
$-r_{N,n}c(\cU)/\norm{\widehat s_\kappa}_\infty$, which dominates $\eta_{N,n}$ because $\eta_{N,n}/r_{N,n}\to0$. Hence the estimated set shrinks to $\cC_0$. Conditional process convergence on the random shrinking set and continuity of the positive-part supremum give size control. Under a fixed violation, the statistic diverges linearly in $r_{N,n}$.
\end{proof}

\begin{remark}[Applications]
For first-order stochastic dominance in the direction
$F^{\cf}\le F^{\fact}$, take $G_0=F^{\cf}-F^{\fact}$. Non-positive tail-effect and quantile-effect hypotheses use the same theorem after replacing the effect surface and derivative loading.
\end{remark}

\section{Numerical approximation and index discretisation}

Suppose the implemented target uses spatial truncation $K$ and numerical time step $\delta$, and let
\begin{equation}\label{eq:eps}
 \norm{\widehat T^{K,\delta}-\widehat T}_\infty
 \le\varepsilon_{K,\delta}.
\end{equation}

\begin{theorem}[Undersmoothed numerical transfer]\label{thm:undersmooth}
The bands and tests above retain their limits if
\[
 r_{N,n}\varepsilon_{K,\delta}\to0.
\]
Thus the sufficient requirements are
\[
 \bn\varepsilon_{K,\delta}\to0
 \quad\text{for fast targets},
 \qquad
 \an\varepsilon_{K,\delta}\to0
 \quad\text{for slow targets}.
\]
Uniform consistency of the bootstrap process additionally requires the corresponding numerical first-derivative errors to vanish uniformly.
\end{theorem}

\begin{theorem}[Bias-enlarged bands]\label{thm:bias-aware}
Suppose a deterministic or high-probability certified bound $B_{K,\delta}$ satisfies
\[
 \norm{T^{K,\delta}_0-T_0}_\infty\le B_{K,\delta}.
\]
Enlarge either raw or studentised half-width by $B_{K,\delta}$. Then the resulting band has asymptotic simultaneous coverage at least $1-\alpha$ without requiring $r_{N,n}B_{K,\delta}\to0$.
\end{theorem}

\begin{proof}
If the statistical band contains the approximate target $T_0^{K,\delta}$, the triangle inequality shows that the enlarged band contains $T_0$. The statistical event has probability tending to $1-\alpha$.
\end{proof}

Let $\cG_m$ be a finite $h_m$-net of the implemented index set.

\begin{proposition}[Grid approximation]\label{prop:grid}
For every continuous function $f$,
\[
 \left|
 \sup_{\tau\in\cT}f(\tau)
 -
 \max_{\tau\in\cG_m}f(\tau)
 \right|
 \le\omega_f(h_m),
\]
where $\omega_f$ is its modulus of continuity. Hence:
\begin{enumerate}[label=(\roman*)]
 \item bootstrap approximation of a continuous-process supremum requires $h_m\to0$;
 \item continuum coverage obtained only from grid endpoints requires either
 \[
 r_{N,n}\omega_{T_0}(h_m)\to0
 \]
 or an explicit interpolation allowance $\omega_{T_0}(h_m)$;
 \item a finite scientific dictionary of weights is an exact special case;
 \item basis truncation of the compact weight ellipsoid must have target error negligible at the relevant fast or slow rate, unless a certified tail allowance is added.
\end{enumerate}
\end{proposition}

\section{Implementation algorithm}

For one data set:
\begin{enumerate}[label=\arabic*.]
 \item Compute the exact conditional MLE $(\thetahat,\sigmahat)$.
 \item Compute the modal score blocks $S_{k,N,n}(\vhat)$.
 \item Scale them by $\Dmat$ and form $\Jhat$.
 \item Draw Gaussian multipliers and obtain $H_{N,n}^*$ from \eqref{eq:bootstrap-score}, or draw directly from $\mathcal N(0,\Jhat^{-1})$.
 \item Evaluate every target derivative surface at the estimated parameter.
 \item Form fast and slow bootstrap target processes using the same draw $H_{N,n}^*$.
 \item Compute raw, floor-studentised, or joint critical values.
 \item For one-sided tests, estimate the contact set using \eqref{eq:Chat} before computing critical values.
 \item Add a numerical-bias allowance when the Galerkin/time approximation is not undersmoothed.
 \item Refine intervention, threshold, and functional grids until critical values and band widths stabilise.
\end{enumerate}

\section{Monte Carlo validation}

The exact-OU experiment uses
\[
 (N,n,K)\in
 \{(5,375,2),(8,960,3),(12,2160,4)\},
 \qquad
 \lambda_N\Delta_n=\frac1{15}.
\]
Each design has 250 Monte Carlo replications and 499 bootstrap draws. Target surfaces use three intervention values, five thresholds, and the first $K$ constant-in-time modal weights. The threshold grid includes zero, where the first-order CDF derivative vanishes.

Table~\ref{tab:coverage} reports the analytic Gaussian multiplier. Monte Carlo standard errors are approximately $0.011$--$0.015$.

\begin{table}[ht]
\centering
\small
\begin{tabular}{@{}cccccccc@{}}
\toprule
$N$ & $n$ & $K$ & Fast raw & Fast stud. & Slow raw & Slow stud. & Nominal\\
\midrule
5  & 375  & 2 & 0.956 & 0.952 & 0.952 & 0.952 & 0.950\\
8  & 960  & 3 & 0.944 & 0.940 & 0.952 & 0.948 & 0.950\\
12 & 2160 & 4 & 0.960 & 0.952 & 0.964 & 0.956 & 0.950\\
\bottomrule
\end{tabular}
\caption{Simultaneous coverage of analytic Gaussian multiplier bands.}
\label{tab:coverage}
\end{table}

The mode-score multiplier is mildly conservative in the smallest design: its slow-effect coverage is $0.980$ at $N=5$, falling to $0.968$ at $N=12$. This reflects covariance estimation from only $N$ independent blocks. The plug-in LAN covariance is therefore the recommended implementation in the positive core. The block multiplier remains important because it verifies the influence construction and generalises to extensions where the limiting covariance is not available in closed form.

\begin{figure}[ht]
\centering
\includegraphics[width=.78\textwidth]{wp5_bootstrap_coverage.png}
\caption{Simultaneous coverage across the three designs.}
\label{fig:coverage}
\end{figure}

\begin{figure}[ht]
\centering
\includegraphics[width=.78\textwidth]{wp5_bootstrap_band_width.png}
\caption{Contraction of average full band width. Fast CDF bands shrink at the $\bn$ scale; slow effect bands shrink at the $\an$ scale.}
\label{fig:width}
\end{figure}

\section{Completion statement and next package}

WP5 meets its completion gate:
\begin{enumerate}[label=(\roman*)]
 \item the mode-block multiplier is feasible and conditionally valid;
 \item simultaneous CDF, quantile, tail, and effect bands are available;
 \item derivative zeros are handled without invalid division;
 \item equality and integrated tests have feasible critical values;
 \item one-sided tests use an estimated contact set;
 \item grid and numerical approximation errors are separated from statistical error;
 \item nominal simultaneous coverage is reproduced across several $(N,n,K)$ designs.
\end{enumerate}

The next package is WP6, not WP7. WP6 must derive semigroup and law derivatives under operator perturbations and define the causal response geometry. WP7 can then combine those operator bounds with the statistical, Galerkin, and temporal errors into the paper's total-error theorem.

\begin{thebibliography}{9}

\bibitem{AndrewsSoares2010}
D. W. K. Andrews and G. Soares.
\newblock Inference for parameters defined by moment inequalities using generalized moment selection.
\newblock \emph{Econometrica}, 78(1):119--157, 2010.

\bibitem{FangSantos2019}
Z. Fang and A. Santos.
\newblock Inference on directionally differentiable functions.
\newblock \emph{Review of Economic Studies}, 86(1):377--412, 2019.

\bibitem{LintonSongWhang2010}
O. Linton, K. Song, and Y.-J. Whang.
\newblock An improved bootstrap test of stochastic dominance.
\newblock \emph{Journal of Econometrics}, 154(2):186--202, 2010.

\bibitem{vdVW1996}
A. W. van der Vaart and J. A. Wellner.
\newblock \emph{Weak Convergence and Empirical Processes}.
\newblock Springer, 1996.

\end{thebibliography}

\end{document}
